%=============================================================================
% PSI-FIELD COMPUTING: LOGIC, MEMORY, AND NEURAL PROCESSING
% IN SCALAR REFRACTIVE GEOMETRY
%
% Patent #10: psi-Field Computing and Logic Devices
% Author: Gary Thomas Alcock
% Density Field Dynamics Research Programme
%=============================================================================

\documentclass[12pt,a4paper,twoside]{article}

%--- Packages ----------------------------------------------------------------
\usepackage[utf8]{inputenc}
\usepackage[T1]{fontenc}
\usepackage{amsmath,amssymb,amsfonts,amsthm}
\usepackage{mathtools}
\usepackage{physics}
\usepackage{siunitx}
\usepackage{graphicx}
\usepackage{float}
\usepackage{booktabs}
\usepackage{array}
\usepackage{multirow}
\usepackage{enumitem}
\usepackage[colorlinks=true,linkcolor=blue!60!black,citecolor=green!50!black,urlcolor=blue!70!black]{hyperref}
\usepackage[margin=2.5cm]{geometry}
\usepackage{fancyhdr}
\usepackage{tcolorbox}
\usepackage{tikz}
\usetikzlibrary{arrows.meta,calc,patterns,decorations.pathmorphing,positioning,shapes.geometric}
\usepackage{pgfplots}
\pgfplotsset{compat=1.18}
\usepackage{caption}
\usepackage{subcaption}
\usepackage{appendix}
\usepackage{xcolor}

%--- Theorem environments ----------------------------------------------------
\newtheorem{theorem}{Theorem}[section]
\newtheorem{proposition}[theorem]{Proposition}
\newtheorem{lemma}[theorem]{Lemma}
\newtheorem{corollary}[theorem]{Corollary}
\newtheorem{definition}[theorem]{Definition}
\newtheorem{remark}[theorem]{Remark}

%--- Custom commands ---------------------------------------------------------
\newcommand{\psifield}{\psi}
\newcommand{\DFD}{DFD}
\newcommand{\kB}{k_{\mathrm{B}}}
\newcommand{\Epsi}{E_\psi}
\newcommand{\Landauer}{E_{\mathrm{L}}}
\newcommand{\Qpsi}{Q_\psi}
\DeclareMathOperator{\sgn}{sgn}
\DeclareMathOperator{\sech}{sech}

%--- Headers -----------------------------------------------------------------
\pagestyle{fancy}
\fancyhf{}
\fancyhead[LE]{\small\textit{\psifield-Field Computing}}
\fancyhead[RO]{\small\textit{G.\,T.\,Alcock}}
\fancyfoot[C]{\thepage}
\renewcommand{\headrulewidth}{0.4pt}

%=============================================================================
\begin{document}

\begin{center}
{\LARGE\bfseries $\psifield$-Field Computing: Logic, Memory, and Neural\\[4pt]
Processing in Scalar Refractive Geometry}

\vspace{1em}

{\large Gary Thomas Alcock}\\[4pt]
{\small Density Field Dynamics Research Programme}\\[2pt]
{\small\texttt{densityfielddynamics.org}}

\vspace{1em}

{\small\today}

\vspace{1.5em}

\begin{minipage}{0.92\textwidth}
\small\textbf{Abstract.}
We develop a complete theoretical framework for computation using the scalar refractive field $\psi$ of Density Field Dynamics (\DFD) as the physical substrate. In \DFD{}, the single scalar field $\psi(\mathbf{x},t)$ defines a refractive index $n = e^\psi$ governing all electromagnetic propagation. We show that this field naturally supports: (i)~bistable potential minima encoding binary logic states, with switching energies that can be driven below the Landauer limit $\kB T \ln 2$ per operation due to the inherent reversibility of the field dynamics; (ii)~coupled-cavity architectures implementing universal logic gates through phase-controlled interference of $\psi$ modes; (iii)~metastable standing-wave configurations providing non-volatile memory; and (iv)~weighted gradient couplings enabling neuromorphic and analog processing. We derive the information-theoretic capacity of $\psi$-field channels, establish energy bounds for logic operations, design complete gate sets (NOT, AND, OR, NAND, XOR), and compare performance projections against CMOS, photonic, quantum, and neuromorphic architectures. The geometric synchronization inherent in cavity-coupled $\psi$ modes eliminates the need for external clocking, while the absence of charge transport removes Joule heating as a fundamental dissipation channel. We present a detailed engineering design for a proof-of-concept $\psi$-logic demonstrator using coupled superconducting radio-frequency (SRF) cavities, and discuss applications to AI acceleration, quantum simulation, and radiation-hardened computing.
\end{minipage}
\end{center}

\vspace{1em}

%=============================================================================
\tableofcontents
\newpage

%=============================================================================
\section{Introduction}\label{sec:intro}
%=============================================================================

The semiconductor industry has sustained exponential performance scaling for over five decades, but fundamental physical limits are now constraining further progress. Transistor gate lengths approaching the atomic scale (${\sim}2$\,nm) encounter quantum tunneling leakage, while power dissipation at densities exceeding $100\,\si{W/cm^2}$ presents insurmountable thermal management challenges. These barriers have motivated exploration of radically different computational substrates: photonic circuits, quantum processors, neuromorphic chips, and various analog computing schemes.

Simultaneously, the thermodynamic foundations of computation have received renewed attention. Landauer's principle establishes that any logically irreversible operation---one that erases information---must dissipate at least $\kB T \ln 2 \approx 2.87 \times 10^{-21}\,\si{J}$ per bit at room temperature. Modern CMOS transistors dissipate roughly $10^4$--$10^5$ times this limit. While reversible computing architectures in principle circumvent Landauer's bound for non-erasing operations, practical implementations in electronic hardware remain challenging due to parasitic resistances, capacitances, and the fundamentally dissipative nature of charge transport through resistive media.

This paper introduces a fundamentally different approach: computation performed directly within the scalar refractive field $\psi$ of Density Field Dynamics (\DFD). In \DFD{}, all electromagnetic propagation occurs through an optical metric defined by a single scalar field:
\begin{equation}
n(\mathbf{x},t) = e^{\psi(\mathbf{x},t)},
\label{eq:refindex}
\end{equation}
where $\psi$ is the fundamental gravitational degree of freedom on flat Minkowski spacetime. The field energy is
\begin{equation}
E_\psi = \int |\nabla\psi|^2\,d^3x,
\label{eq:field_energy}
\end{equation}
in the linearized (high-gradient) regime, and the field equation for quasi-static configurations is the nonlinear Poisson equation
\begin{equation}
\nabla \cdot \left[\mu\!\left(\frac{|\nabla\psi|}{a_*}\right)\nabla\psi\right] = \frac{4\pi G}{c^2}(\rho - \bar{\rho}),
\label{eq:field_eq}
\end{equation}
where $\mu(x)$ is the interpolating function and $a_* = 2a_0/c^2$ is the characteristic gradient scale.

The key insight is that $\psi$-field dynamics in engineered cavities are:
\begin{enumerate}[label=(\roman*)]
\item \textbf{Energy-conserving}: The field equation derives from a Lagrangian with a convex energy functional, so lossless dynamics are the natural state; dissipation must be \emph{engineered in}, not suppressed.
\item \textbf{Wave-supporting}: Perturbations $\delta\psi$ propagate as waves with well-defined modes in bounded cavities, enabling interference-based logic.
\item \textbf{Nonlinear}: The $\mu$-function provides intrinsic nonlinearity for bistable switching without external amplification.
\item \textbf{Geometrically synchronized}: Cavity mode frequencies are set by geometry, providing natural clock signals without external oscillators.
\end{enumerate}

These properties make $\psi$-field cavities a natural platform for reversible, energy-efficient computation.

\subsection{Relation to Existing Computing Paradigms}\label{sec:paradigm_comparison}

The $\psi$-field computing architecture intersects several active research directions while being distinct from each:

\paragraph{Optical computing.}
Photonic circuits manipulate electromagnetic fields to perform logic. However, conventional optical computing faces challenges with cascadability (output intensity degrades through successive gates), fan-out, and the weakness of optical nonlinearities. $\psi$-field computing differs fundamentally: it operates on the \emph{medium} through which light propagates, not on the light itself. The refractive field provides strong effective nonlinearity through its self-interaction potential.

\paragraph{Reversible computing.}
The field of reversible computing, pioneered by Bennett, Fredkin, and Toffoli, seeks to minimize information erasure and thus energy dissipation. Electronic implementations face practical limits from parasitic dissipation in wires and switches. $\psi$-field computing achieves reversibility naturally: the Lagrangian dynamics conserve energy, and information is encoded in field configurations rather than charge states.

\paragraph{Neuromorphic computing.}
Physical neural networks using photonic meshes, memristive crossbars, or spintronic arrays map weighted connections onto physical couplings. $\psi$-field computing extends this concept: weighted gradient couplings between cavity modes directly implement synaptic connections, with coupling strengths set by geometric overlap integrals.

\paragraph{Reservoir computing.}
Reservoir computers exploit the complex dynamics of physical systems (optical cavities, mechanical oscillators, spin networks) for computation. $\psi$-field cavities provide a rich dynamical reservoir: the nonlinear mode interactions and spatial complexity of $\psi$ configurations can be harnessed for temporal pattern recognition.

\subsection{Paper Organization}

Section~\ref{sec:info_theory} develops the information-theoretic foundations of $\psi$-field computation. Section~\ref{sec:bistable} analyzes bistable $\psi$ potentials and binary encoding. Section~\ref{sec:gates} designs a complete set of logic gates using coupled cavities. Section~\ref{sec:memory} addresses memory storage in metastable standing waves. Section~\ref{sec:energy} derives energy dissipation bounds and the sub-Landauer regime. Section~\ref{sec:neural} develops the neuromorphic and analog processing architecture. Section~\ref{sec:comparison} compares $\psi$-field computing with CMOS, quantum, photonic, and other paradigms. Section~\ref{sec:applications} discusses applications, and Section~\ref{sec:demonstrator} presents the engineering design for a proof-of-concept demonstrator.

%=============================================================================
\section{Information Theory of the $\psi$-Field}\label{sec:info_theory}
%=============================================================================

\subsection{Information Content of Field Configurations}\label{sec:info_content}

A $\psi$-field configuration in a cavity of volume $V$ encodes information in the spatial and temporal structure of the field. We quantify this capacity by decomposing $\psi$ into cavity modes.

Consider a rectangular cavity with dimensions $L_x \times L_y \times L_z$. The field admits a mode expansion:
\begin{equation}
\psi(\mathbf{x},t) = \psi_0 + \sum_{\mathbf{n}} q_{\mathbf{n}}(t)\,\phi_{\mathbf{n}}(\mathbf{x}),
\label{eq:mode_expansion}
\end{equation}
where $\phi_{\mathbf{n}}(\mathbf{x})$ are the orthonormal cavity eigenmodes satisfying Dirichlet or Neumann boundary conditions, and $q_{\mathbf{n}}(t)$ are the mode amplitudes. Each mode $\mathbf{n} = (n_x, n_y, n_z)$ has eigenfrequency
\begin{equation}
\omega_{\mathbf{n}} = c_\psi \pi \sqrt{\frac{n_x^2}{L_x^2} + \frac{n_y^2}{L_y^2} + \frac{n_z^2}{L_z^2}},
\label{eq:mode_freq}
\end{equation}
where $c_\psi$ is the $\psi$-wave propagation speed in the cavity medium.

\begin{definition}[Modal Information Capacity]
The information capacity of a $\psi$-field cavity with $N_{\rm mode}$ accessible modes, each supporting $M$ distinguishable amplitude levels, is
\begin{equation}
\mathcal{C}_\psi = N_{\rm mode} \log_2 M \quad \text{bits}.
\label{eq:info_capacity}
\end{equation}
The number of distinguishable levels per mode is determined by the signal-to-noise ratio:
\begin{equation}
M = 1 + \frac{q_{\rm max}}{\delta q_{\rm min}},
\label{eq:levels}
\end{equation}
where $q_{\rm max}$ is the maximum mode amplitude (set by nonlinear saturation) and $\delta q_{\rm min}$ is the minimum resolvable amplitude (set by thermal noise or quantum fluctuations).
\end{definition}

\subsection{Thermal Noise Floor}\label{sec:noise}

The thermal energy per mode in equilibrium is $\frac{1}{2}\kB T$, giving an RMS amplitude fluctuation:
\begin{equation}
\langle q_{\mathbf{n}}^2 \rangle_{\rm thermal} = \frac{\kB T}{M_\psi \omega_{\mathbf{n}}^2},
\label{eq:thermal_noise}
\end{equation}
where $M_\psi$ is the effective mass of the $\psi$ mode (related to the energy functional through $E_{\mathbf{n}} = \frac{1}{2}M_\psi \omega_{\mathbf{n}}^2 q_{\mathbf{n}}^2$).

The quantum noise floor sets a harder limit:
\begin{equation}
\langle q_{\mathbf{n}}^2 \rangle_{\rm quantum} = \frac{\hbar}{2M_\psi \omega_{\mathbf{n}}},
\label{eq:quantum_noise}
\end{equation}
corresponding to zero-point fluctuations.

\subsection{Channel Capacity for $\psi$-Field Communication}\label{sec:channel}

For a $\psi$-field channel of bandwidth $B$ connecting two coupled cavities, the Shannon capacity is
\begin{equation}
C = B \log_2\!\left(1 + \mathrm{SNR}_\psi\right),
\label{eq:shannon}
\end{equation}
where the signal-to-noise ratio for $\psi$-mode signaling is
\begin{equation}
\mathrm{SNR}_\psi = \frac{P_{\rm signal}}{P_{\rm noise}} = \frac{M_\psi \omega^2 q_{\rm sig}^2}{\kB T \cdot B}.
\label{eq:snr}
\end{equation}
For a cavity with mode frequency $\omega/2\pi = 1\,\si{GHz}$, quality factor $Q = 10^{10}$ (achievable in SRF cavities), and operating temperature $T = 1.5\,\si{K}$:
\begin{equation}
\mathrm{SNR}_\psi \sim \frac{\hbar\omega Q}{\kB T} \approx \frac{(6.6\times 10^{-25})(10^{10})}{(1.4\times 10^{-23})(1.5)} \approx 300,
\label{eq:snr_estimate}
\end{equation}
giving a capacity per mode of $C \approx 8$\,bits per mode period.

\subsection{Information Density}\label{sec:info_density}

The volumetric information density of a $\psi$-field memory is bounded by the number of independent modes per unit volume up to a cutoff frequency $\omega_{\rm max}$:
\begin{equation}
\rho_{\rm info} = \frac{1}{6\pi^2}\left(\frac{\omega_{\rm max}}{c_\psi}\right)^3 \log_2 M \quad \text{bits/m}^3.
\label{eq:info_density}
\end{equation}
For $c_\psi \sim c$ and $\omega_{\rm max}/2\pi = 10\,\si{GHz}$:
\begin{equation}
\rho_{\rm info} \sim 10^6 \times \log_2 M \quad \text{bits/m}^3.
\end{equation}
This is sparse compared to solid-state memory ($\sim 10^{25}\,\text{bits/m}^3$), but the $\psi$-field memory operates with fundamentally different energy-speed-reversibility trade-offs (see Section~\ref{sec:energy}).

%=============================================================================
\section{Bistable $\psi$-Potentials and Binary Encoding}\label{sec:bistable}
%=============================================================================

\subsection{Double-Well Potential for Binary States}\label{sec:double_well}

Binary logic requires two distinguishable, stable states. In $\psi$-field computing, these are provided by the minima of an effective double-well potential for a cavity mode amplitude $q$.

The nonlinear $\psi$-field self-interaction, combined with boundary conditions and external biasing, generates an effective potential for each mode:
\begin{equation}
V_{\rm eff}(q) = \frac{1}{2}M_\psi \omega_0^2 q^2 + \frac{\beta}{4}q^4 - h\,q + V_{\rm bias}(q),
\label{eq:double_well}
\end{equation}
where:
\begin{itemize}
\item $M_\psi \omega_0^2$ is the linear restoring force (from the cavity geometry).
\item $\beta$ is the quartic nonlinearity coefficient arising from the nonlinear kinetic function $W(y)$ in the DFD action.
\item $h$ is an external bias field (controllable).
\item $V_{\rm bias}$ represents coupling to adjacent cavities.
\end{itemize}

For the binary computing application, we engineer the effective potential to have the form:
\begin{equation}
V_{\rm eff}(q) = \frac{\beta}{4}(q^2 - q_0^2)^2,
\label{eq:symmetric_dw}
\end{equation}
which has minima at $q = \pm q_0$ with barrier height
\begin{equation}
\Delta V = \frac{\beta q_0^4}{4}.
\label{eq:barrier}
\end{equation}

The logical assignment is:
\begin{equation}
|0\rangle_\psi \equiv q = -q_0, \qquad |1\rangle_\psi \equiv q = +q_0.
\label{eq:logical_states}
\end{equation}

\subsection{Switching Dynamics}\label{sec:switching}

Transitions between logic states occur via controlled barrier traversal. The switching can be driven by:

\paragraph{(a) Bias-driven switching.}
Applying a bias field $h$ tilts the double well, lowering the barrier on one side:
\begin{equation}
\Delta V_{\rm eff}(h) = \Delta V\left(1 - \frac{h}{h_c}\right)^2, \qquad h_c = \frac{2\beta q_0^3}{3\sqrt{3}},
\label{eq:biased_barrier}
\end{equation}
where $h_c$ is the critical bias at which the metastable minimum vanishes. For $h > h_c$, switching is deterministic and rapid.

\paragraph{(b) Resonant driving.}
Driving the mode at its natural frequency $\omega_0$ builds amplitude until the barrier is surmounted:
\begin{equation}
q(t) = q_{\rm init} + \frac{F_0}{M_\psi}\frac{\sin(\omega_0 t)}{\gamma_\psi},
\end{equation}
where $F_0$ is the driving force amplitude and $\gamma_\psi$ is the mode damping rate.

\paragraph{(c) Parametric switching.}
Modulating the cavity boundary conditions at $2\omega_0$ parametrically amplifies the mode until it crosses the barrier. This is the most energy-efficient switching mechanism, as the driving energy can be recovered.

\subsection{Switching Speed and Energy}\label{sec:switch_speed}

The fundamental switching time is bounded by the oscillation period of the mode:
\begin{equation}
\tau_{\rm switch} \geq \frac{\pi}{\omega_0} = \frac{1}{2f_0}.
\label{eq:switch_time}
\end{equation}
For a mode at $f_0 = 1\,\si{GHz}$, $\tau_{\rm switch} \geq 0.5\,\si{ns}$, comparable to current CMOS switching speeds.

The minimum energy required for a reliable (error-free) switching event is set by the barrier height, which must exceed thermal fluctuations:
\begin{equation}
\Delta V \geq \alpha_{\rm rel}\,\kB T,
\label{eq:reliability}
\end{equation}
where $\alpha_{\rm rel} \approx 40$--$60$ for error rates below $10^{-15}$ (comparable to CMOS reliability targets). At room temperature, this gives $\Delta V \geq 10^{-19}\,\si{J}$.

\subsection{Phase-Space Representation}\label{sec:phase_space}

The state of a $\psi$ mode is fully characterized by the phase-space coordinates $(q, p)$ where $p = M_\psi \dot{q}$. The logical states are well-separated regions in phase space:
\begin{equation}
\Gamma_0 = \{(q,p) : q < 0,\; E(q,p) < \Delta V\}, \quad \Gamma_1 = \{(q,p) : q > 0,\; E(q,p) < \Delta V\},
\label{eq:phase_regions}
\end{equation}
where $E(q,p) = p^2/(2M_\psi) + V_{\rm eff}(q)$. The Liouville theorem guarantees that the phase-space volume is conserved under Hamiltonian evolution, providing the foundation for reversible logic.

%=============================================================================
\section{Logic Gate Design}\label{sec:gates}
%=============================================================================

\subsection{The Coupled-Cavity Architecture}\label{sec:coupled_cavity}

Logic gates are implemented by coupling two or more $\psi$-field cavities through apertures or waveguide sections. The coupling allows energy and information exchange between cavity modes.

Consider two cavities (A and B) coupled through a connecting channel. The Hamiltonian for the coupled system is:
\begin{equation}
H = \frac{p_A^2}{2M_A} + V_A(q_A) + \frac{p_B^2}{2M_B} + V_B(q_B) + J(q_A, q_B),
\label{eq:coupled_ham}
\end{equation}
where $J(q_A, q_B)$ is the coupling potential. For linear coupling:
\begin{equation}
J_{\rm lin} = \kappa \, q_A q_B,
\label{eq:linear_coupling}
\end{equation}
with coupling strength $\kappa$ determined by the geometric overlap integral:
\begin{equation}
\kappa = c_\psi^2 \int_{\rm aperture} \nabla\phi_A \cdot \nabla\phi_B \, dS,
\label{eq:coupling_strength}
\end{equation}
where $\phi_{A,B}$ are the mode functions and the integral is over the coupling aperture.

\subsection{$\psi$-NOT Gate}\label{sec:not_gate}

The NOT gate inverts the logical state: $|0\rangle_\psi \to |1\rangle_\psi$ and $|1\rangle_\psi \to |0\rangle_\psi$.

\paragraph{Implementation.} A single cavity with an inverting bias pulse. The input state $q_{\rm in} = \pm q_0$ is driven to $q_{\rm out} = \mp q_0$ by a calibrated bias pulse $h(t)$:
\begin{equation}
h(t) = h_0 \operatorname{sech}\!\left(\frac{t - t_0}{\tau_p}\right) \operatorname{sgn}(q_{\rm in}),
\end{equation}
where $h_0 > h_c$ ensures deterministic switching, and $\operatorname{sech}$ provides an adiabatic pulse shape that minimizes ringing.

\paragraph{Energy.} The NOT gate is logically reversible (bijective), so it can in principle operate with zero net energy dissipation. The bias pulse energy is deposited and then recovered:
\begin{equation}
E_{\rm NOT} = \oint h(t)\,\dot{q}\,dt = 0 \quad \text{(ideal reversible)}.
\end{equation}
In practice, finite $Q$ leads to dissipation $E_{\rm NOT} \sim \Delta V / Q_\psi$.

\subsection{$\psi$-NAND Gate (Universal Gate)}\label{sec:nand_gate}

The NAND gate is functionally complete---any Boolean function can be constructed from NAND gates alone.

\paragraph{Architecture.} Three coupled cavities: two input cavities (A, B) and one output cavity (C). The coupling is designed so that the output state depends on the product of input states:

\begin{equation}
H_{\rm NAND} = H_A + H_B + H_C + J_{AC}(q_A, q_C) + J_{BC}(q_B, q_C) + J_{\rm NL}(q_A, q_B, q_C),
\label{eq:nand_ham}
\end{equation}
where the nonlinear three-body coupling is:
\begin{equation}
J_{\rm NL} = \lambda_3 \, q_A q_B q_C.
\label{eq:three_body}
\end{equation}

\paragraph{Truth table.}
The nonlinear coupling $J_{\rm NL}$ and the bias on cavity C are tuned so that:
\begin{center}
\begin{tabular}{cc|c|l}
\toprule
$q_A$ & $q_B$ & $q_C$ (output) & Logic \\
\midrule
$-q_0$ ($|0\rangle$) & $-q_0$ ($|0\rangle$) & $+q_0$ ($|1\rangle$) & $\overline{0 \cdot 0} = 1$ \\
$-q_0$ ($|0\rangle$) & $+q_0$ ($|1\rangle$) & $+q_0$ ($|1\rangle$) & $\overline{0 \cdot 1} = 1$ \\
$+q_0$ ($|1\rangle$) & $-q_0$ ($|0\rangle$) & $+q_0$ ($|1\rangle$) & $\overline{1 \cdot 0} = 1$ \\
$+q_0$ ($|1\rangle$) & $+q_0$ ($|1\rangle$) & $-q_0$ ($|0\rangle$) & $\overline{1 \cdot 1} = 0$ \\
\bottomrule
\end{tabular}
\end{center}

\paragraph{Mechanism.} The three-body coupling $\lambda_3 q_A q_B q_C$ generates an effective bias on $q_C$ proportional to $q_A q_B$. When both inputs are in the $|1\rangle$ state ($q_A = q_B = +q_0$), the product $q_A q_B = +q_0^2$ biases C toward $-q_0$ (logical 0). For all other input combinations, $q_A q_B \leq 0$ or $= q_0^2$ but with opposite sign contribution, biasing C toward $+q_0$ (logical 1).

\subsection{$\psi$-AND Gate}\label{sec:and_gate}

\paragraph{Implementation.} NAND followed by NOT, or equivalently a direct two-input coupling with opposite bias polarity:
\begin{equation}
J_{\rm AND} = -\lambda_3\,q_A q_B q_C + h_{\rm offset}\,q_C,
\end{equation}
where $h_{\rm offset}$ is tuned so that only when both inputs are $|1\rangle$ does the output settle to $|1\rangle$.

\subsection{$\psi$-OR Gate}\label{sec:or_gate}

\paragraph{Implementation.} Using De Morgan's law: $A \vee B = \overline{\bar{A} \cdot \bar{B}}$. Alternatively, direct implementation via sum coupling:
\begin{equation}
J_{\rm OR} = \kappa_A q_A q_C + \kappa_B q_B q_C + \frac{\lambda_2}{2}(q_A + q_B)^2 q_C,
\end{equation}
with bias chosen so that the output is $|1\rangle$ whenever at least one input is $|1\rangle$.

\subsection{$\psi$-XOR Gate}\label{sec:xor_gate}

The XOR gate is particularly natural in wave-based computing, as it corresponds to destructive vs.\ constructive interference.

\paragraph{Implementation.} Two input cavities couple to the output through paths with a $\pi$ relative phase shift for same-sign inputs:
\begin{equation}
J_{\rm XOR} = \kappa(q_A - q_B)^2 q_C - \frac{\kappa}{q_0^2}\,q_A q_B\,q_C.
\label{eq:xor_coupling}
\end{equation}
When $q_A = q_B$ (both 0 or both 1), the coupling vanishes and C relaxes to $|0\rangle$. When $q_A \neq q_B$, the coupling drives C to $|1\rangle$.

\subsection{Reversible Fredkin and Toffoli Gates}\label{sec:reversible_gates}

For fully reversible computation, we implement the Fredkin (controlled-SWAP) and Toffoli (controlled-controlled-NOT) gates.

\paragraph{Toffoli gate.} Three-cavity implementation with inputs $q_A$, $q_B$ (controls) and $q_C$ (target):
\begin{equation}
H_{\rm Toffoli} = H_A + H_B + H_C + \Lambda_4\,q_A q_B(q_C^2 - q_0^2) + h_T\,q_A q_B q_C,
\end{equation}
where the quartic term $\Lambda_4 q_A q_B(q_C^2 - q_0^2)$ modulates the barrier height of cavity C conditioned on the states of A and B. Only when both controls are $|1\rangle$ ($q_A = q_B = +q_0$) is the barrier lowered, allowing the bias term $h_T$ to flip C.

\paragraph{Fredkin gate.} Controlled swap of two target cavities B and C, controlled by A:
\begin{equation}
H_{\rm Fredkin} = H_A + H_B + H_C + \kappa_F\,q_A(q_B - q_C)^2,
\end{equation}
where the coupling $\kappa_F q_A(q_B - q_C)^2$ only activates when $q_A = +q_0$ (control = 1), creating an energy penalty for $q_B \neq q_C$ that drives B and C toward their swapped values.

\subsection{Fan-Out and Signal Restoration}\label{sec:fanout}

\paragraph{Fan-out.} A single $\psi$-mode output drives multiple downstream cavities through independent coupling apertures. The geometric isolation of different apertures ensures minimal back-action:
\begin{equation}
\kappa_{\rm back} \sim \kappa_{\rm forward} \cdot e^{-L/\ell_{\rm decay}},
\end{equation}
where $L$ is the separation between output apertures and $\ell_{\rm decay}$ is the evanescent decay length in the coupling waveguide.

\paragraph{Signal restoration.} The bistable potential automatically restores degraded signals. Any mode amplitude within the basin of attraction of a logical state ($|q - q_0| < q_0$, approximately) will relax to the precise logical value $q_0$ via the double-well dynamics. This self-restoration is analogous to CMOS signal restoration but achieved through field dynamics rather than transistor gain.

\subsection{Gate Timing and Synchronization}\label{sec:timing}

\paragraph{Geometric clock.} The mode frequencies of identical cavities are identical by construction, providing a natural clock. A chain of $N$ coupled cavities with uniform geometry oscillates at frequencies
\begin{equation}
\omega_k = \omega_0\sqrt{1 + 2\frac{\kappa}{\omega_0^2}\cos\frac{k\pi}{N+1}}, \qquad k = 1, \ldots, N,
\end{equation}
forming a tight band around $\omega_0$. The bandwidth $\Delta\omega \approx 2\kappa/\omega_0$ sets the synchronization tolerance.

\paragraph{Pipeline stages.} Sequential logic is implemented by coupling cavity chains where each chain represents one pipeline stage. The propagation delay through a coupling section provides the pipeline clock:
\begin{equation}
t_{\rm stage} = \frac{L_{\rm coupling}}{c_\psi} + \tau_{\rm settle},
\end{equation}
where $\tau_{\rm settle} \sim Q/\omega_0$ is the settling time for the output to reach its logical state.

%=============================================================================
\section{Memory: Metastable Standing-Wave Configurations}\label{sec:memory}
%=============================================================================

\subsection{Static Memory (SRAM Analogue)}\label{sec:sram}

A $\psi$-field static memory cell is a single bistable cavity maintaining its state through the double-well potential. The retention time is set by the thermal activation rate over the barrier:
\begin{equation}
\tau_{\rm retain} = \tau_0 \exp\!\left(\frac{\Delta V}{\kB T}\right),
\label{eq:retention}
\end{equation}
where $\tau_0 \sim 1/\omega_0$ is the attempt frequency. For $\Delta V = 50\,\kB T$ at $T = 1.5\,\si{K}$:
\begin{equation}
\tau_{\rm retain} \sim 10^{-9} \times e^{50} \approx 5 \times 10^{12}\,\si{s} \approx 160{,}000\,\text{years}.
\end{equation}
This exceeds any practical requirement, demonstrating that $\psi$-field memory is effectively non-volatile at cryogenic temperatures.

\subsection{Dynamic Memory (DRAM Analogue)}\label{sec:dram}

Higher-density memory can use shallow double wells ($\Delta V \sim 10\,\kB T$) that require periodic refresh. The refresh rate is:
\begin{equation}
f_{\rm refresh} > \frac{1}{\tau_{\rm retain}} = \frac{\omega_0}{2\pi}\,e^{-\Delta V/\kB T}.
\end{equation}
For $\Delta V = 10\,\kB T$ at $1.5\,\si{K}$ with $\omega_0/2\pi = 1\,\si{GHz}$: $f_{\rm refresh} \sim 50\,\si{Hz}$.

\subsection{Multi-Valued Memory}\label{sec:multi_memory}

The $\psi$-field naturally supports multi-valued storage by engineering potentials with more than two minima:
\begin{equation}
V_N(q) = V_0\left[1 - \cos\!\left(\frac{N\pi q}{q_{\rm max}}\right)\right],
\end{equation}
which has $N$ local minima, encoding $\log_2 N$ bits per cavity. The barrier between adjacent minima is $2V_0$. For $N = 4$ (2 bits per cavity), the barrier is halved relative to the binary case, requiring proportionally lower temperatures for equivalent reliability.

\subsection{Standing-Wave Pattern Memory}\label{sec:sw_memory}

A qualitatively different memory architecture exploits the spatial structure of standing-wave patterns in extended cavities. A one-dimensional cavity of length $L$ supports standing-wave patterns:
\begin{equation}
\psi_{\rm mem}(x) = \sum_{n=1}^{N} a_n \sin\!\left(\frac{n\pi x}{L}\right),
\label{eq:sw_pattern}
\end{equation}
where the coefficient vector $\{a_n\}$ encodes the stored data. This is a form of holographic memory, where information is distributed across the spatial extent of the cavity.

The storage capacity of a pattern memory of length $L$ with maximum mode number $N$ is:
\begin{equation}
\mathcal{C}_{\rm pattern} = N \log_2 M_{\rm eff}\quad \text{bits},
\end{equation}
where $M_{\rm eff}$ is the number of distinguishable amplitude levels per mode.

\paragraph{Read-out.} The stored pattern is read by coupling a probe $\psi$-mode into the cavity and measuring the reflected/transmitted amplitude spectrum. The probe acquires phase shifts proportional to the mode amplitudes:
\begin{equation}
\Delta\phi_n = \frac{2\omega_{\rm probe}}{c_\psi}\,\frac{\partial n}{\partial \psi}\bigg|_{\psi_0}\! a_n L_{\rm eff},
\end{equation}
enabling non-destructive read-out.

\subsection{Write Operations}\label{sec:write}

Writing to a $\psi$-field memory cell requires driving the mode from one basin of attraction to the other:

\paragraph{Method 1: Direct bias.} Apply an external $\psi$-field gradient via a control electrode (piezoelectric actuator modifying boundary conditions):
\begin{equation}
h_{\rm write}(t) = h_c\,\Pi\!\left(\frac{t - t_w}{\tau_w}\right),
\end{equation}
where $\Pi$ is a rectangular pulse of duration $\tau_w > \pi/\omega_0$.

\paragraph{Method 2: Resonant injection.} Inject a $\psi$-wave through the coupling aperture at the mode frequency, building amplitude until the barrier is crossed. This is slower ($\tau_w \sim Q/\omega_0$ cycles) but more energy-efficient.

%=============================================================================
\section{Energy Dissipation and Sub-Landauer Operation}\label{sec:energy}
%=============================================================================

\subsection{Landauer's Principle Revisited}\label{sec:landauer}

Landauer's principle states that erasing one bit of information in a system at temperature $T$ requires a minimum energy dissipation:
\begin{equation}
\Landauer = \kB T \ln 2 \approx 2.87 \times 10^{-21}\,\si{J} \quad \text{at } T = 300\,\si{K}.
\label{eq:landauer}
\end{equation}

This bound applies to \emph{logically irreversible} operations---those that map multiple input states to a single output state (e.g., AND, OR, ERASE). Logically \emph{reversible} operations (bijective maps like NOT, SWAP, Toffoli) are not subject to this bound and can in principle operate at zero dissipation.

At cryogenic temperatures ($T = 1.5\,\si{K}$), the Landauer limit drops dramatically:
\begin{equation}
\Landauer(1.5\,\si{K}) = 1.4 \times 10^{-23}\,\si{J} \approx 0.09\,\si{meV}.
\end{equation}

\subsection{Energy Budget of $\psi$-Field Logic}\label{sec:energy_budget}

The total energy per logic operation in a $\psi$-field gate has three contributions:
\begin{equation}
E_{\rm op} = E_{\rm switch} + E_{\rm dissipate} + E_{\rm signal},
\label{eq:energy_total}
\end{equation}
where:

\paragraph{(a) Switching energy $E_{\rm switch}$.}
This is the energy temporarily stored in the field during the switching transient. For the double-well potential~\eqref{eq:symmetric_dw}:
\begin{equation}
E_{\rm switch} = \Delta V = \frac{\beta q_0^4}{4}.
\end{equation}
In a reversible operation, this energy is fully recovered after the transient, so $E_{\rm switch}$ is borrowed, not dissipated.

\paragraph{(b) Dissipated energy $E_{\rm dissipate}$.}
The fraction of switching energy lost to the cavity walls during the switching transient:
\begin{equation}
E_{\rm dissipate} = E_{\rm switch} \cdot \frac{\pi}{\omega_0 \tau_{\rm switch}} \cdot \frac{1}{Q_\psi} = \frac{\pi \Delta V}{Q_\psi},
\label{eq:dissipation}
\end{equation}
where $Q_\psi$ is the quality factor of the $\psi$-mode and we have assumed the switching takes one half-period ($\tau_{\rm switch} = \pi/\omega_0$).

\paragraph{(c) Signal transmission energy $E_{\rm signal}$.}
Energy radiated into the coupling channels to communicate the output to downstream gates:
\begin{equation}
E_{\rm signal} = \frac{\kappa^2 q_0^2}{M_\psi \omega_0^2} \cdot N_{\rm fan},
\end{equation}
where $N_{\rm fan}$ is the fan-out count. This energy is absorbed by downstream cavities and used for their switching---it is not wasted.

\subsection{The Sub-Landauer Regime}\label{sec:sub_landauer}

\begin{theorem}[Sub-Landauer Operation of $\psi$-Logic]
\label{thm:sub_landauer}
A reversible $\psi$-field logic gate operating at temperature $T$ with cavity quality factor $Q_\psi$ dissipates energy per operation
\begin{equation}
E_{\rm op}^{\rm rev} = \frac{\pi\alpha_{\rm rel}\,\kB T}{Q_\psi},
\label{eq:sub_landauer}
\end{equation}
where $\alpha_{\rm rel} \geq 40$ is the reliability factor. This is below the Landauer limit when
\begin{equation}
Q_\psi > \frac{\pi\alpha_{\rm rel}}{\ln 2} \approx 182.
\label{eq:Q_threshold}
\end{equation}
\end{theorem}

\begin{proof}
For a reversible operation (NOT, SWAP, Toffoli), no information is erased, so the Landauer bound does not apply. The only dissipation is the wall loss during switching:
\begin{equation}
E_{\rm dissipate} = \frac{\pi\Delta V}{Q_\psi} = \frac{\pi\alpha_{\rm rel}\,\kB T}{Q_\psi},
\end{equation}
using $\Delta V = \alpha_{\rm rel}\,\kB T$ for reliable switching. This is less than $\kB T \ln 2$ when $Q_\psi > \pi\alpha_{\rm rel}/\ln 2$.
\end{proof}

For SRF cavities with $Q_\psi \sim 10^{10}$:
\begin{equation}
E_{\rm op}^{\rm rev}(1.5\,\si{K}) = \frac{\pi \times 50 \times (1.4 \times 10^{-23}) \times 1.5}{10^{10}} \approx 3.3 \times 10^{-31}\,\si{J},
\label{eq:srf_dissipation}
\end{equation}
which is roughly \textbf{ten billion times below the room-temperature Landauer limit} and $10^{8}$ times below the cryogenic Landauer limit.

\begin{remark}
Even for logically irreversible gates (AND, OR), the dissipation per operation is $E_{\rm op}^{\rm irrev} = \kB T \ln 2 + \pi\alpha_{\rm rel}\kB T/Q_\psi$. With high-$Q$ cavities, the irreversible dissipation is dominated by the Landauer term, not wall losses. This motivates implementing all computation in reversible form (Toffoli/Fredkin gates with ancilla management).
\end{remark}

\subsection{Comparison of Energy Dissipation}\label{sec:energy_comparison}

\begin{table}[H]
\centering
\caption{Energy per logic operation across computing architectures.}
\label{tab:energy}
\begin{tabular}{lcc}
\toprule
\textbf{Architecture} & \textbf{Energy/op (J)} & \textbf{Ratio to Landauer (300\,K)} \\
\midrule
Modern CMOS (7\,nm) & $\sim 10^{-17}$ & $\sim 3{,}500$ \\
Advanced CMOS (2\,nm, projected) & $\sim 10^{-18}$ & $\sim 350$ \\
Superconducting logic (SFQ) & $\sim 10^{-19}$ & $\sim 35$ \\
Photonic logic (ring resonators) & $\sim 10^{-18}$ & $\sim 350$ \\
Landauer limit (300\,K) & $2.87\times 10^{-21}$ & $1$ \\
Landauer limit (1.5\,K) & $1.4\times 10^{-23}$ & $0.005$ \\
$\psi$-field logic (SRF, reversible) & $\sim 3\times 10^{-31}$ & $\sim 10^{-10}$ \\
\bottomrule
\end{tabular}
\end{table}

\subsection{Power Density Analysis}\label{sec:power_density}

For a $\psi$-field processor operating at clock frequency $f_{\rm clk} = 1\,\si{GHz}$ with $N_{\rm gates} = 10^6$ simultaneously switching gates:
\begin{equation}
P_{\rm total} = N_{\rm gates} \cdot E_{\rm op} \cdot f_{\rm clk} = 10^6 \times 3 \times 10^{-31} \times 10^9 = 3 \times 10^{-16}\,\si{W}.
\end{equation}
This is negligible compared to the cryogenic cooling power required ($\sim 1$--$100\,\si{W}$), demonstrating that for SRF-based $\psi$-field computing, the energy budget is dominated by the cryogenics, not the computation itself.

%=============================================================================
\section{Neuromorphic and Analog Processing}\label{sec:neural}
%=============================================================================

\subsection{Weighted Gradient Couplings as Synapses}\label{sec:synapses}

The coupling between $\psi$-field cavities provides a natural implementation of synaptic connections. The coupling strength $\kappa_{ij}$ between cavities $i$ and $j$ plays the role of the synaptic weight:
\begin{equation}
\kappa_{ij} = c_\psi^2 \int_{\rm aperture_{ij}} \nabla\phi_i \cdot \nabla\phi_j\,dS.
\label{eq:synaptic_weight}
\end{equation}

The weight is controlled by:
\begin{enumerate}
\item \textbf{Aperture geometry}: The size and shape of the coupling aperture.
\item \textbf{Aperture medium}: Inserting a material with $\psi$-dependent properties (e.g., a tunable dielectric) in the coupling channel.
\item \textbf{Intermediate mode}: Placing a mediating cavity in the coupling path with a controllable detuning $\delta\omega$.
\end{enumerate}

The effective weight through an intermediary with detuning $\delta\omega$ is:
\begin{equation}
\kappa_{ij}^{\rm eff} = \frac{\kappa_{i,\rm med}\,\kappa_{\rm med,j}}{\delta\omega^2 + \gamma_{\rm med}^2},
\label{eq:mediated_weight}
\end{equation}
providing a continuously tunable weight by varying $\delta\omega$.

\subsection{Activation Functions from Nonlinear Mode Response}\label{sec:activation}

The double-well potential provides a natural sigmoid activation function. The steady-state response of a bistable cavity to an input signal $s$ is:
\begin{equation}
q_{\rm out}(s) = q_0\,\tanh\!\left(\frac{s}{s_0}\right),
\label{eq:activation}
\end{equation}
where $s_0 = 2\Delta V / q_0$ is the sensitivity scale. This is precisely the hyperbolic tangent activation function used in neural networks.

Other activation functions can be engineered by shaping the potential:
\begin{itemize}
\item \textbf{ReLU analogue}: Asymmetric double well with one flat branch: $V(q) = \beta(q - q_0)^2\Theta(q - q_c)$.
\item \textbf{Softplus}: Logarithmic potential: $V(q) = V_0 \ln(1 + e^{q/q_c})$.
\item \textbf{Step function}: Deep double well with sharp barrier (high $\beta$).
\end{itemize}

\subsection{Neural Layer Architecture}\label{sec:neural_layer}

A single neural network layer is implemented as:
\begin{equation}
q_j^{(l+1)} = \sigma\!\left(\sum_i \kappa_{ij}\,q_i^{(l)} + b_j\right),
\label{eq:neural_layer}
\end{equation}
where $q_i^{(l)}$ are the mode amplitudes of layer-$l$ cavities, $\kappa_{ij}$ are the coupling weights, $b_j$ is a bias (set by an external field gradient on cavity $j$), and $\sigma$ is the activation function from the bistable response.

\paragraph{Physical implementation.} Each layer is a row of $N_l$ cavities. The coupling matrix $\kappa_{ij}$ is implemented by the network of coupling apertures between layers. The matrix-vector multiplication $\sum_i \kappa_{ij} q_i^{(l)}$ happens \emph{at the speed of wave propagation}---it is an inherently parallel, analog operation.

\subsection{Learning and Weight Update}\label{sec:learning}

Weight updates in the $\psi$-field neural network require modifying coupling strengths. Two mechanisms are available:

\paragraph{(a) Geometric reconfiguration.} Mechanically adjustable apertures (using piezoelectric actuators) change the geometric overlap integral~\eqref{eq:synaptic_weight}. This is slow ($\sim \si{ms}$) but provides persistent weight storage.

\paragraph{(b) Parametric tuning.} Modulating the intermediary cavity detuning $\delta\omega$ in Eq.~\eqref{eq:mediated_weight} electronically adjusts weights on nanosecond timescales. This enables rapid weight updates for online learning.

\paragraph{(c) $\psi$-field plasticity.} If the $\psi$ field exhibits hysteresis in certain nonlinear regimes, the coupling strength can be modified by the history of mode amplitudes---a form of physical Hebbian learning:
\begin{equation}
\frac{d\kappa_{ij}}{dt} = \eta\,q_i(t)\,q_j(t) - \lambda\,\kappa_{ij},
\label{eq:hebbian}
\end{equation}
where $\eta$ is the learning rate and $\lambda$ is a decay constant preventing runaway growth.

\subsection{Reservoir Computing Mode}\label{sec:reservoir}

A large $\psi$-field cavity with many interacting modes naturally functions as a reservoir computer. The high-dimensional nonlinear dynamics of the mode interactions map temporal input signals into a rich feature space.

\paragraph{Architecture.}
\begin{enumerate}
\item \textbf{Input}: Inject a time-varying $\psi$-wave into a large, irregularly shaped cavity.
\item \textbf{Reservoir}: The cavity modes interact nonlinearly via the $W(y)$ potential, generating a complex dynamical response.
\item \textbf{Readout}: Sample the mode amplitudes at multiple spatial points. These form the feature vector $\mathbf{x}(t)$.
\item \textbf{Training}: Only the readout weights $\mathbf{w}$ are trained (linear regression), giving output $y(t) = \mathbf{w} \cdot \mathbf{x}(t)$.
\end{enumerate}

The computational power of the reservoir depends on its \emph{fading memory} and \emph{separation property}, both of which are naturally provided by the damped nonlinear $\psi$-field dynamics.

\subsection{Comparison with Existing Neuromorphic Architectures}\label{sec:neuro_comparison}

\begin{table}[H]
\centering
\caption{Neuromorphic architecture comparison.}
\label{tab:neuro}
\begin{tabular}{lcccl}
\toprule
\textbf{Platform} & \textbf{Energy/MAC} & \textbf{Speed} & \textbf{Weight bits} & \textbf{Mechanism} \\
\midrule
Intel Loihi 2 & $\sim 10^{-14}\,\si{J}$ & $\sim 1\,\si{GHz}$ & 8 & Digital spiking \\
IBM TrueNorth & $\sim 10^{-13}\,\si{J}$ & $\sim 1\,\si{kHz}$ & 1--4 & Digital spiking \\
Photonic mesh & $\sim 10^{-17}\,\si{J}$ & $\sim 10\,\si{GHz}$ & $>8$ & MZI arrays \\
Memristive array & $\sim 10^{-15}\,\si{J}$ & $\sim 100\,\si{MHz}$ & 4--6 & $R$ states \\
$\psi$-field neural & $\sim 10^{-31}\,\si{J}$ & $\sim 1\,\si{GHz}$ & analog & Gradient coupling \\
\bottomrule
\end{tabular}
\end{table}

The energy per multiply-accumulate (MAC) operation in $\psi$-field neural processing is dramatically lower than all existing platforms, though this advantage is offset by the cryogenic operating requirement.

%=============================================================================
\section{Comparison with Computing Paradigms}\label{sec:comparison}
%=============================================================================

\subsection{$\psi$-Field vs.\ CMOS}\label{sec:vs_cmos}

\begin{table}[H]
\centering
\caption{CMOS vs.\ $\psi$-field computing comparison.}
\label{tab:vs_cmos}
\begin{tabular}{lll}
\toprule
\textbf{Property} & \textbf{CMOS} & \textbf{$\psi$-Field} \\
\midrule
State variable & Charge (electrons) & Field amplitude ($q_{\mathbf{n}}$) \\
Switching mechanism & Transistor on/off & Barrier traversal \\
Energy dissipation & Joule heating & Wall losses (no charge transport) \\
Reversibility & Requires special circuits & Natural (Hamiltonian dynamics) \\
Clock & External oscillator & Geometric (cavity modes) \\
Fan-out & Charge redistribution & Wave coupling (no loading) \\
Noise immunity & Voltage margins & Basin of attraction \\
Scalability limit & Tunneling, heat & Cavity miniaturization, $Q$ \\
Operating temperature & Room temp & Cryogenic (for high $Q$) \\
\bottomrule
\end{tabular}
\end{table}

\subsection{$\psi$-Field vs.\ Quantum Computing}\label{sec:vs_quantum}

$\psi$-field computing is \emph{not} quantum computing---it uses classical field configurations, not quantum superpositions. However, the two architectures share SRF cavity technology and cryogenic infrastructure.

Key differences:
\begin{itemize}
\item \textbf{Decoherence}: Quantum computers require maintaining fragile quantum superpositions; $\psi$-field computers use robust classical states (bistable potentials) that are naturally error-resistant.
\item \textbf{Error correction}: Quantum error correction requires massive overhead ($10^3$--$10^4$ physical qubits per logical qubit). $\psi$-field self-restoration provides automatic error correction at no overhead.
\item \textbf{Universality}: Quantum computers provide exponential speedup for certain problems (factoring, simulation). $\psi$-field computers provide polynomial speedup through massive parallelism and energy efficiency.
\item \textbf{Programming model}: Quantum algorithms require careful unitary circuit design. $\psi$-field programs map naturally to neural network and analog computing models.
\end{itemize}

\subsection{$\psi$-Field vs.\ Photonic Computing}\label{sec:vs_photonic}

Photonic computing manipulates electromagnetic fields directly; $\psi$-field computing manipulates the \emph{medium} those fields propagate through. This distinction has important consequences:

\begin{itemize}
\item \textbf{Nonlinearity}: Optical nonlinearities are weak ($\chi^{(3)} \sim 10^{-22}\,\si{m^2/V^2}$ in silica), requiring high intensities. The $\psi$-field self-interaction is inherently strong (the $W(y)$ function is order-unity nonlinear).
\item \textbf{Cascadability}: Photonic gate outputs degrade through successive stages due to loss and nonlinearity constraints. $\psi$-field bistable restoration regenerates clean logical states at each gate.
\item \textbf{Memory}: Photonic systems struggle with compact optical memory. $\psi$-field bistable potentials provide dense, non-volatile memory.
\end{itemize}

\subsection{$\psi$-Field vs.\ Superconducting Logic}\label{sec:vs_sfq}

Single-flux-quantum (SFQ) logic is the closest existing technology to $\psi$-field computing, sharing cryogenic operation and electromagnetic field-based state variables.

\begin{itemize}
\item \textbf{State variable}: SFQ uses magnetic flux quanta ($\Phi_0 = h/2e$) in Josephson junction loops. $\psi$-field uses mode amplitudes in cavities.
\item \textbf{Speed}: SFQ operates at $\sim 100\,\si{GHz}$ (limited by junction capacitance). $\psi$-field operates at $\sim 1$--$10\,\si{GHz}$ (limited by cavity mode frequency).
\item \textbf{Energy}: SFQ dissipates $\sim 10^{-19}\,\si{J}$/op. $\psi$-field dissipates $\sim 10^{-31}\,\si{J}$/op (reversible).
\item \textbf{Density}: SFQ achieves moderate integration. $\psi$-field faces challenges with cavity size ($\sim \si{cm}$ for GHz modes).
\end{itemize}

%=============================================================================
\section{Applications}\label{sec:applications}
%=============================================================================

\subsection{AI Acceleration}\label{sec:ai}

The neuromorphic architecture of Section~\ref{sec:neural} naturally accelerates neural network inference and training:

\paragraph{Matrix multiplication.} The coupling matrix $\kappa_{ij}$ between cavity layers performs matrix-vector multiplication at the speed of wave propagation. For an $N \times N$ weight matrix, this is an $O(1)$ time operation (independent of $N$), compared to $O(N^2)$ for digital and $O(N)$ for systolic array implementations.

\paragraph{Energy efficiency.} The energy per MAC of $\sim 10^{-31}\,\si{J}$ translates to:
\begin{equation}
\text{TOPS/W} = \frac{10^{12}}{10^{-31}} = 10^{43}\,\si{TOPS/W},
\end{equation}
though this theoretical peak is unrealistically high due to cryogenic overhead. Including a realistic cryogenic system ($\sim 1\,\si{kW}$ wall power for $\sim 1\,\si{W}$ cooling at $1.5\,\si{K}$):
\begin{equation}
\text{TOPS/W}_{\rm system} \sim \frac{N_{\rm MAC} \cdot f}{P_{\rm cryo}} \sim \frac{10^6 \times 10^9}{10^3} = 10^{12}\,\si{TOPS/W},
\end{equation}
which would exceed current GPU efficiency ($\sim 10^3\,\si{TOPS/W}$) by nine orders of magnitude for sufficiently large systems.

\paragraph{Transformer acceleration.} The self-attention mechanism in transformer architectures requires computing $\mathbf{Q}\mathbf{K}^T$ (query-key products) and softmax normalization. A $\psi$-field implementation maps:
\begin{itemize}
\item Query/key/value projections $\to$ coupling matrices between cavity layers.
\item Attention scores $\to$ interference amplitudes between query and key cavities.
\item Softmax $\to$ nonlinear cavity response with competition (mutual inhibition).
\end{itemize}

\subsection{Quantum Simulation}\label{sec:quantum_sim}

While $\psi$-field computing is classical, it can efficiently simulate certain quantum systems:

\paragraph{Lattice gauge theories.} The $\psi$-field on a discrete cavity lattice is governed by:
\begin{equation}
H_{\rm lattice} = \sum_i V(q_i) + \sum_{\langle ij \rangle} \kappa_{ij}(q_i - q_j)^2,
\end{equation}
which has the same form as a lattice scalar field theory. By tuning the potential $V$ and coupling $\kappa$, different effective field theories can be simulated.

\paragraph{Ising model.} The bistable cavities with nearest-neighbor coupling directly implement the Ising Hamiltonian:
\begin{equation}
H_{\rm Ising} = -\sum_{\langle ij \rangle} J_{ij} \sigma_i \sigma_j - \sum_i h_i \sigma_i,
\end{equation}
where $\sigma_i = q_i/q_0 = \pm 1$ and $J_{ij} = \kappa_{ij} q_0^2$. Adiabatic cooling of the $\psi$-field system naturally finds the ground state---an analog of quantum annealing.

\subsection{Radiation-Hardened Computing}\label{sec:rad_hard}

$\psi$-field computing is inherently radiation-tolerant:
\begin{enumerate}
\item \textbf{No charge collection}: Unlike semiconductor devices, where ionizing radiation generates electron-hole pairs that corrupt stored charge, $\psi$-field states are field configurations that cannot be disrupted by individual particle interactions.
\item \textbf{Distributed storage}: Information is stored in the global mode pattern of a cavity, not in a localized charge packet. A localized perturbation (particle hit) couples weakly to the mode.
\item \textbf{Self-restoration}: The bistable potential automatically corrects small perturbations to the field configuration.
\end{enumerate}

The single-event upset (SEU) cross-section for a $\psi$-field memory cell is estimated as:
\begin{equation}
\sigma_{\rm SEU} \sim \sigma_{\rm geom} \cdot \left(\frac{E_{\rm deposit}}{E_{\rm barrier}}\right)^{N_{\rm mode}},
\end{equation}
where $N_{\rm mode}$ is the number of modes involved in storing the bit. For $N_{\rm mode} > 1$, the SEU rate is exponentially suppressed.

\subsection{Signal Processing}\label{sec:signal}

The wave nature of $\psi$-field computation provides natural implementations of signal processing operations:
\begin{itemize}
\item \textbf{Fourier transform}: Coupling a signal into a multi-mode cavity automatically decomposes it into frequency components (the cavity modes).
\item \textbf{Convolution}: Passing a signal through a shaped $\psi$-field medium implements spatial convolution.
\item \textbf{Matched filtering}: A cavity pre-loaded with a template pattern correlates with incoming signals through mode overlap.
\end{itemize}

%=============================================================================
\section{Engineering Design: Proof-of-Concept Demonstrator}\label{sec:demonstrator}
%=============================================================================

\subsection{Design Overview}\label{sec:design_overview}

The proof-of-concept demonstrator implements a single $\psi$-NAND gate plus one memory cell, sufficient to demonstrate:
\begin{enumerate}
\item Bistable $\psi$-field states in an SRF cavity.
\item Controlled switching between states.
\item Logic gate operation (NAND truth table).
\item Non-volatile $\psi$-field memory.
\end{enumerate}

\subsection{Cavity Specifications}\label{sec:cavity_specs}

\begin{table}[H]
\centering
\caption{SRF cavity specifications for the $\psi$-logic demonstrator.}
\label{tab:cavity_specs}
\begin{tabular}{lll}
\toprule
\textbf{Parameter} & \textbf{Value} & \textbf{Rationale} \\
\midrule
Cavity type & Niobium elliptical & Standard SRF technology \\
Fundamental frequency & $1.3\,\si{GHz}$ & TESLA/ILC standard \\
Quality factor $Q_0$ & $> 10^{10}$ & Demonstrated in DESY/JLAB cavities \\
Operating temperature & $1.8\,\si{K}$ (superfluid He) & Below Nb $T_c = 9.2\,\si{K}$ \\
Cavity length & $23\,\si{cm}$ (single cell) & $\lambda/2$ at 1.3\,GHz \\
Cavity diameter & $18.7\,\si{cm}$ & TESLA geometry \\
Aperture diameter & $7\,\si{cm}$ & Standard iris \\
$E_{\rm acc}$ (max) & $25\,\si{MV/m}$ & Conservative operating point \\
$R_s$ (surface resistance) & $< 10\,\si{n\Omega}$ & State-of-the-art Nb preparation \\
\bottomrule
\end{tabular}
\end{table}

\subsection{NAND Gate Implementation}\label{sec:nand_impl}

\paragraph{Physical layout.} Three single-cell SRF cavities (A, B, C) connected by coupling waveguide sections:

\begin{center}
\begin{tikzpicture}[scale=0.8, every node/.style={font=\small}]
  % Cavity A
  \draw[thick, fill=blue!10] (0,2) ellipse (1.2 and 0.6);
  \node at (0,2) {Cavity A};
  \node[above] at (0,2.8) {\textbf{Input 1}};

  % Cavity B
  \draw[thick, fill=blue!10] (0,-2) ellipse (1.2 and 0.6);
  \node at (0,-2) {Cavity B};
  \node[below] at (0,-2.8) {\textbf{Input 2}};

  % Coupling waveguides
  \draw[thick, -Stealth] (1.2,1.8) -- (3.5,0.3);
  \draw[thick, -Stealth] (1.2,-1.8) -- (3.5,-0.3);
  \node[right] at (2.2,1.2) {$\kappa_{AC}$};
  \node[right] at (2.2,-1.2) {$\kappa_{BC}$};

  % Nonlinear coupler
  \draw[thick, fill=red!10] (3.5,0) circle (0.3);
  \node[above] at (3.5,0.5) {NL};

  % Cavity C
  \draw[thick, fill=green!10] (5.5,0) ellipse (1.2 and 0.6);
  \node at (5.5,0) {Cavity C};
  \node[above] at (5.5,0.8) {\textbf{Output}};

  % Arrow from NL to C
  \draw[thick, -Stealth] (3.8,0) -- (4.3,0);
\end{tikzpicture}
\end{center}

\paragraph{Nonlinear coupler.} The three-body coupling $\lambda_3 q_A q_B q_C$ is implemented by a nonlinear element at the junction of the three waveguides. Two approaches:

\begin{enumerate}
\item \textbf{Parametric mixing via Josephson junction}: A Josephson junction in the coupling region provides the required $\cos\phi$ nonlinearity. The junction is biased to operate in its nonlinear regime, where the current-phase relation $I = I_c \sin\phi$ generates harmonic mixing of the input signals.

\item \textbf{Ferrite-based nonlinearity}: A ferrite sphere placed at the coupling junction provides magnetic nonlinearity. The nonlinear permeability $\mu(H) = \mu_0(1 + \chi H + \chi_3 H^3 + \cdots)$ generates the required third-order mixing.
\end{enumerate}

\subsection{Memory Cell Implementation}\label{sec:mem_impl}

\paragraph{Physical layout.} A single SRF cavity with:
\begin{itemize}
\item Adjustable end-wall position (piezoelectric actuator) for bias control.
\item A Josephson junction on the cavity wall providing the quartic nonlinearity for the double-well potential.
\item Weakly coupled probe antenna for non-destructive read-out.
\end{itemize}

The effective potential is:
\begin{equation}
V_{\rm eff}(q) = \frac{1}{2}M_\psi\omega_0^2 q^2 - \frac{E_J}{2}\cos\!\left(\frac{2eq}{\hbar}\right) \approx \frac{1}{2}(M_\psi\omega_0^2 - E_J k_J^2)q^2 + \frac{E_J k_J^4}{24}q^4,
\end{equation}
where $E_J$ is the Josephson energy and $k_J = 2e/\hbar$. When $E_J k_J^2 > M_\psi\omega_0^2$, the quadratic term is negative and the potential becomes a double well.

\subsection{Read-Out System}\label{sec:readout}

\paragraph{Phase-sensitive detection.} A low-power probe tone is injected into the cavity through a weakly coupled port. The reflected probe acquires a phase shift that depends on the mode amplitude:
\begin{equation}
\Delta\phi_{\rm probe} = \frac{2\omega_{\rm probe}}{c}\frac{dn}{d\psi}\bigg|_{\psi_0}\,q\,L_{\rm eff} = \pm\frac{2\omega_{\rm probe}}{c}e^{\psi_0}\,q_0\,L_{\rm eff}.
\end{equation}

The sign of $\Delta\phi_{\rm probe}$ determines the logical state. A homodyne detection system with a local oscillator achieves single-shot readout in $\tau_{\rm read} \sim 1/\Delta f_{\rm probe} \sim 10\,\si{ns}$.

\subsection{Control and Bias System}\label{sec:control}

\begin{itemize}
\item \textbf{Bias generation}: Low-noise current sources driving magnetic coils near the cavity produce controllable $\psi$-field gradients (via EM--$\psi$ coupling with parameter $\lambda$).
\item \textbf{Parametric drive}: RF synthesizers modulating cavity boundary conditions at $2\omega_0$ for parametric switching.
\item \textbf{Timing}: A master oscillator phase-locked to the cavity fundamental provides the system clock.
\item \textbf{Temperature control}: Superfluid helium bath at $1.8\,\si{K}$ with $\pm 1\,\si{mK}$ stability.
\end{itemize}

\subsection{Expected Experimental Signatures}\label{sec:signatures}

The demonstrator should produce the following observable signatures:

\begin{enumerate}
\item \textbf{Bistability}: The cavity mode amplitude exhibits hysteresis as a function of applied bias, with a clear jump between two stable values at the critical bias $h_c$.

\item \textbf{Switching}: Application of a bias pulse $h > h_c$ causes deterministic transition from one stable state to the other, observable as a sign change in the probe phase $\Delta\phi_{\rm probe}$.

\item \textbf{NAND truth table}: The output cavity C settles to the correct logical state for all four input combinations, with a settling time $\tau_{\rm settle} < 100\,\si{ns}$.

\item \textbf{Memory retention}: After writing a state and removing the bias, the memory cell retains its state for $> 10^6$ seconds (limited by experimental patience, not physics).

\item \textbf{Energy measurement}: Calorimetric measurement of heat dissipated per switching event, expected to be below $10^{-23}\,\si{J}$ at $1.8\,\si{K}$.
\end{enumerate}

\subsection{Demonstrator Bill of Materials}\label{sec:bom}

\begin{table}[H]
\centering
\caption{Key components for the $\psi$-logic demonstrator.}
\label{tab:bom}
\begin{tabular}{llp{5cm}}
\toprule
\textbf{Component} & \textbf{Qty} & \textbf{Notes} \\
\midrule
Nb SRF cavities (1.3\,GHz, single-cell) & 4 & TESLA geometry; 3 for NAND + 1 for memory \\
Coupling waveguide sections & 3 & WR-650 or custom Nb waveguide \\
Josephson junction assemblies & 2 & For NL coupler + memory double-well \\
Piezoelectric tuners & 4 & Cavity frequency fine-tuning \\
Cryostat (superfluid He) & 1 & $\geq 50\,\si{L}$ capacity, $1.8\,\si{K}$ \\
RF synthesizers (1.3\,GHz) & 3 & Phase-coherent, $< -140\,\si{dBc/Hz}$ \\
Low-noise amplifiers (cryogenic) & 4 & HEMT, $T_n < 5\,\si{K}$ \\
Homodyne detection electronics & 4 & IQ mixers + ADC \\
Magnetic shielding & 1 set & $\mu$-metal + superconducting shield \\
FPGA control system & 1 & Timing, sequencing, data acquisition \\
\bottomrule
\end{tabular}
\end{table}

\subsection{Scaling Roadmap}\label{sec:scaling}

\begin{enumerate}
\item \textbf{Phase 1} (Demonstrator): Single NAND gate + memory cell. Validates bistability, switching, logic. Timeline: 18--24 months.
\item \textbf{Phase 2} (4-bit processor): 16 coupled cavities implementing a 4-bit ALU. Demonstrates cascaded logic, pipelining, arithmetic. Timeline: 3--4 years.
\item \textbf{Phase 3} (Neuromorphic array): $32 \times 32$ cavity array implementing a neural network layer. Demonstrates parallel MAC operations, learning. Timeline: 5--7 years.
\item \textbf{Phase 4} (Integrated system): Multi-layer neural processor with $> 10^4$ cavities. Miniaturized cavities (higher frequencies), integrated cryogenics. Timeline: 8--12 years.
\end{enumerate}

\paragraph{Miniaturization path.} Moving to higher frequencies enables smaller cavities: at $10\,\si{GHz}$, a cavity is $\sim 1.5\,\si{cm}$; at $100\,\si{GHz}$, $\sim 1.5\,\si{mm}$. However, $Q$ typically decreases at higher frequencies, requiring advances in surface treatment and material science.

%=============================================================================
\section{Discussion}\label{sec:discussion}
%=============================================================================

\subsection{Advantages of $\psi$-Field Computing}\label{sec:advantages}

\begin{enumerate}
\item \textbf{Fundamental reversibility}: The Lagrangian field dynamics are energy-conserving by construction. Dissipation is a perturbation (wall losses), not a fundamental feature.

\item \textbf{No charge transport}: The absence of moving charges eliminates Joule heating, electromigration, and shot noise---the three dominant dissipation and reliability mechanisms in electronic computing.

\item \textbf{Geometric synchronization}: Cavity mode frequencies provide natural clock signals, eliminating clock distribution networks and their associated power consumption.

\item \textbf{Self-restoring signals}: Bistable potentials automatically correct signal degradation, providing built-in error correction without redundancy overhead.

\item \textbf{Intrinsic nonlinearity}: The $W(y)$ function provides strong nonlinearity at the field level, unlike photonic systems that require high intensities or external nonlinear materials.

\item \textbf{Radiation hardness}: Field-based state storage is inherently resistant to single-event upsets from ionizing radiation.

\item \textbf{Analog/digital duality}: The same cavity can operate as a binary (bistable) or analog (multi-level/continuous) element, providing flexible computing modes.
\end{enumerate}

\subsection{Challenges and Limitations}\label{sec:challenges}

\begin{enumerate}
\item \textbf{Cryogenic requirement}: High $Q$ factors require superconducting cavities operating at $T < 4\,\si{K}$. The cryogenic infrastructure dominates system cost and power. Room-temperature $\psi$-field computing would require discovering high-$Q$ cavity materials operable at $300\,\si{K}$.

\item \textbf{Physical size}: At $1.3\,\si{GHz}$, each cavity is ${\sim}20\,\si{cm}$ long. A 4-bit processor would occupy ${\sim}1\,\si{m}^3$, far larger than electronic equivalents. Miniaturization requires moving to higher frequencies with maintained $Q$.

\item \textbf{Integration density}: The fundamental limit on cavity size is set by the wavelength of the mode, giving a minimum gate pitch of $\lambda/2 \sim c/(2f)$. At $100\,\si{GHz}$, this is $1.5\,\si{mm}$---comparable to early integrated circuits but far larger than modern CMOS.

\item \textbf{$\psi$-field coupling verification}: The EM--$\psi$ coupling parameter $\lambda$ has not yet been experimentally measured. If $|\lambda - 1|$ is very small, the interaction channel for biasing may be weak, requiring alternative control mechanisms.

\item \textbf{Fabrication precision}: The coupling strengths $\kappa_{ij}$ must be controlled to high precision for accurate neural network weights. This requires sub-micron machining tolerances on cavity geometries.
\end{enumerate}

\subsection{Implications for the DFD Programme}\label{sec:implications}

The $\psi$-field computing architecture provides a technology pathway for testing DFD predictions:

\begin{enumerate}
\item A functioning $\psi$-logic gate would constitute direct evidence for the EM--$\psi$ coupling and the physical reality of the scalar refractive field.
\item The bistability measurement would verify the nonlinear self-interaction potential of $\psi$.
\item Precision measurement of switching energies would constrain the $\psi$-field parameters ($M_\psi$, $\beta$, $\lambda$) to unprecedented accuracy.
\end{enumerate}

%=============================================================================
\section{Conclusions}\label{sec:conclusions}
%=============================================================================

We have developed a comprehensive theoretical framework for computing using the scalar refractive field $\psi$ of Density Field Dynamics. The key results are:

\begin{enumerate}
\item \textbf{Binary logic}: Bistable $\psi$-field potentials provide robust binary state encoding, with switching via bias-driven, resonant, or parametric mechanisms.

\item \textbf{Universal gate set}: Coupled-cavity architectures implement NOT, AND, OR, NAND, XOR, and reversible (Toffoli, Fredkin) gates through controlled interference and nonlinear coupling.

\item \textbf{Sub-Landauer energy}: Reversible $\psi$-logic in SRF cavities ($Q \sim 10^{10}$) dissipates $\sim 10^{-31}\,\si{J}$/operation at $1.8\,\si{K}$, roughly $10^{10}$ times below the room-temperature Landauer limit.

\item \textbf{Non-volatile memory}: Metastable standing-wave configurations in double-well potentials provide retention times exceeding $10^{12}\,\si{s}$ at cryogenic temperatures, with multi-valued and pattern-based variants.

\item \textbf{Neuromorphic processing}: Weighted gradient couplings between cavities naturally implement synaptic connections, with continuous weight tunability and analog activation functions arising from the nonlinear cavity response.

\item \textbf{Demonstrator design}: A proof-of-concept system using four SRF cavities at $1.3\,\si{GHz}$ can validate the core principles within existing accelerator technology capabilities.
\end{enumerate}

The $\psi$-field computing paradigm occupies a unique position in the landscape of post-CMOS computing: it combines the reversibility advantages of adiabatic/ballistic logic, the wave-based processing of photonic computing, the neuromorphic flexibility of physical neural networks, and the energy efficiency enabled by superconducting technology---all unified within a single scalar field framework.

The primary challenges are the cryogenic operating requirement and the relatively large physical size of current cavity technology. However, both of these mirror the early-stage challenges faced by quantum computing, which has progressed from laboratory curiosity to practical deployment through sustained engineering effort. The $\psi$-field platform may follow a similar trajectory, particularly for applications where extreme energy efficiency (AI acceleration, high-performance computing) or radiation hardness (space systems, nuclear environments) justify the cryogenic infrastructure.

Most fundamentally, a working $\psi$-field logic device would constitute direct experimental evidence for the physical reality of the DFD scalar refractive field and its interaction with electromagnetic systems---making $\psi$-field computing both a technological application and a precision test of DFD.

%=============================================================================
\begin{appendices}

\section{Derivation of the Effective Double-Well Potential}\label{app:double_well}

We derive the effective double-well potential for a single $\psi$-mode in a cavity with a Josephson junction nonlinearity.

\paragraph{Setup.} Consider a cavity mode with amplitude $q$ and frequency $\omega_0$, coupled to a Josephson junction with energy $E_J$ and phase $\varphi = k_J q$ where $k_J = 2e/\hbar$.

The total energy is:
\begin{equation}
E(q, \dot{q}) = \frac{1}{2}M_\psi\dot{q}^2 + \frac{1}{2}M_\psi\omega_0^2 q^2 - E_J\cos(k_J q).
\end{equation}

\paragraph{Expansion.} For $|k_J q| < \pi$, expand the cosine:
\begin{equation}
E_J\cos(k_J q) \approx E_J\left(1 - \frac{k_J^2 q^2}{2} + \frac{k_J^4 q^4}{24}\right).
\end{equation}

The effective potential is:
\begin{equation}
V_{\rm eff}(q) = \frac{1}{2}\underbrace{(M_\psi\omega_0^2 - E_J k_J^2)}_{\equiv -\alpha_{\rm eff}} q^2 + \frac{E_J k_J^4}{24}q^4 + \text{const.}
\end{equation}

\paragraph{Double-well condition.} When $E_J k_J^2 > M_\psi\omega_0^2$ (i.e., $\alpha_{\rm eff} > 0$):
\begin{equation}
V_{\rm eff}(q) = -\frac{\alpha_{\rm eff}}{2}q^2 + \frac{\beta}{4}q^4, \qquad \beta = \frac{E_J k_J^4}{6}.
\end{equation}

The minima are at:
\begin{equation}
q_0 = \pm\sqrt{\frac{\alpha_{\rm eff}}{\beta}} = \pm\sqrt{\frac{6(E_J k_J^2 - M_\psi\omega_0^2)}{E_J k_J^4}},
\end{equation}
and the barrier height is:
\begin{equation}
\Delta V = \frac{\alpha_{\rm eff}^2}{4\beta} = \frac{3(E_J k_J^2 - M_\psi\omega_0^2)^2}{2E_J k_J^4}.
\end{equation}

\section{Reversibility and Entropy Production}\label{app:reversibility}

\paragraph{Liouville's theorem.} The Hamiltonian dynamics of the coupled-cavity system conserve the phase-space volume element $\prod_i dq_i\,dp_i$. By Liouville's theorem, the entropy $S = -\kB \int f \ln f\,d\Gamma$ is constant for a Hamiltonian system.

\paragraph{Entropy production from dissipation.} The cavity wall losses introduce a friction term $-2\gamma_\psi \dot{q}$ in the equation of motion. The entropy production rate per mode is:
\begin{equation}
\dot{S}_{\rm mode} = \frac{2\gamma_\psi M_\psi \langle\dot{q}^2\rangle}{T} = \frac{2\gamma_\psi E_{\rm kinetic}}{T}.
\end{equation}

During a switching event lasting one half-period, the entropy produced is:
\begin{equation}
\Delta S_{\rm switch} = \frac{2\gamma_\psi}{T}\int_0^{\pi/\omega_0} M_\psi\dot{q}^2\,dt = \frac{\pi\gamma_\psi M_\psi\omega_0 q_0^2}{T} = \frac{\pi\Delta V}{\omega_0 Q_\psi T},
\end{equation}
using $\gamma_\psi = \omega_0/(2Q_\psi)$ and $M_\psi\omega_0^2 q_0^2/2 = \Delta V$.

The corresponding energy dissipation is:
\begin{equation}
E_{\rm dissipate} = T\Delta S_{\rm switch} = \frac{\pi\Delta V}{Q_\psi},
\end{equation}
confirming Eq.~\eqref{eq:dissipation}.

\section{Coupling Strength Calculations}\label{app:coupling}

\paragraph{Aperture coupling.} For two identical cavities connected by a circular aperture of radius $a$, the coupling coefficient (Bethe theory) is:
\begin{equation}
\kappa = \frac{8\omega_0}{3\pi}\left(\frac{a}{L}\right)^3,
\end{equation}
where $L$ is the cavity length. For the TESLA cavity parameters ($a = 3.5\,\si{cm}$, $L = 23\,\si{cm}$):
\begin{equation}
\kappa = \frac{8 \times 2\pi \times 1.3 \times 10^9}{3\pi}\left(\frac{0.035}{0.23}\right)^3 \approx 6.5 \times 10^6\,\si{rad/s} \approx 1\,\si{MHz}.
\end{equation}

This gives a coupling-to-linewidth ratio:
\begin{equation}
\frac{\kappa}{\gamma_\psi} = \frac{\kappa Q_\psi}{\omega_0} = \frac{6.5 \times 10^6 \times 10^{10}}{8.2 \times 10^9} \approx 8 \times 10^6,
\end{equation}
firmly in the strong-coupling regime---essential for efficient logic operation.

\section{Noise Analysis and Error Rates}\label{app:noise}

\paragraph{Thermal switching error rate.} The probability of a thermally activated false switching event per attempt time $\tau_0 = 1/\omega_0$ is:
\begin{equation}
P_{\rm error} = e^{-\Delta V / \kB T}.
\end{equation}

For $\Delta V = 50\,\kB T$: $P_{\rm error} = e^{-50} \approx 2 \times 10^{-22}$, far below the target $10^{-15}$.

\paragraph{Quantum tunneling rate.} The WKB tunneling rate through the double-well barrier is:
\begin{equation}
\Gamma_{\rm tunnel} = \omega_0\,e^{-S_{\rm inst}/\hbar},
\end{equation}
where the instanton action is:
\begin{equation}
S_{\rm inst} = \int_{-q_0}^{+q_0}\sqrt{2M_\psi V_{\rm eff}(q)}\,dq = \frac{4}{3}\sqrt{\frac{2M_\psi}{\beta}}\,\alpha_{\rm eff}^{3/2}.
\end{equation}

For macroscopic cavity modes, $S_{\rm inst}/\hbar \gg 1$, making quantum tunneling completely negligible.

\end{appendices}

%=============================================================================
% REFERENCES
%=============================================================================
\begin{thebibliography}{99}

\bibitem{Alcock2024}
G.\,T.\,Alcock, ``Density Field Dynamics: A Complete Unified Theory,'' v3.2 (2024). \texttt{densityfielddynamics.org}

\bibitem{Gordon1923}
W.\,Gordon, ``Zur Lichtfortpflanzung nach der Relativit\"atstheorie,'' \emph{Ann.\ Phys.} \textbf{72}, 421--456 (1923).

\bibitem{Perlick2000}
V.\,Perlick, \emph{Ray Optics, Fermat's Principle, and Applications to General Relativity}, Lecture Notes in Physics \textbf{61} (Springer, 2000).

\bibitem{Landauer1961}
R.\,Landauer, ``Irreversibility and heat generation in the computing process,'' \emph{IBM J.\ Res.\ Dev.} \textbf{5}, 183--191 (1961).

\bibitem{Bennett1973}
C.\,H.\,Bennett, ``Logical reversibility of computation,'' \emph{IBM J.\ Res.\ Dev.} \textbf{17}, 525--532 (1973).

\bibitem{Bennett1982}
C.\,H.\,Bennett, ``The thermodynamics of computation---a review,'' \emph{Int.\ J.\ Theor.\ Phys.} \textbf{21}, 905--940 (1982).

\bibitem{Fredkin1982}
E.\,Fredkin and T.\,Toffoli, ``Conservative logic,'' \emph{Int.\ J.\ Theor.\ Phys.} \textbf{21}, 219--253 (1982).

\bibitem{BekensteinMilgrom1984}
J.\,Bekenstein and M.\,Milgrom, ``Does the missing mass problem signal the breakdown of Newtonian gravity?,'' \emph{Astrophys.\ J.} \textbf{286}, 7--14 (1984).

\bibitem{Berut2012}
A.\,B\'erut \emph{et al.}, ``Experimental verification of Landauer's principle linking information and thermodynamics,'' \emph{Nature} \textbf{483}, 187--189 (2012).

\bibitem{Frank2005}
M.\,P.\,Frank, ``Introduction to reversible computing: motivation, progress, and challenges,'' \emph{Proc.\ CF'05}, 385--390 (ACM, 2005).

\bibitem{Shen2017}
Y.\,Shen \emph{et al.}, ``Deep learning with coherent nanophotonic circuits,'' \emph{Nat.\ Photon.} \textbf{11}, 441--446 (2017).

\bibitem{Wetzstein2020}
G.\,Wetzstein \emph{et al.}, ``Inference in artificial intelligence with deep optics and photonics,'' \emph{Nature} \textbf{588}, 39--47 (2020).

\bibitem{Wright2022}
L.\,G.\,Wright \emph{et al.}, ``Deep physical neural networks trained with backpropagation,'' \emph{Nature} \textbf{601}, 549--555 (2022).

\bibitem{Momeni2023}
A.\,Momeni \emph{et al.}, ``Backpropagation-free training of deep physical neural networks,'' \emph{Science} \textbf{382}, 1297--1303 (2023).

\bibitem{Shastri2021}
B.\,J.\,Shastri \emph{et al.}, ``Photonics for artificial intelligence and neuromorphic computing,'' \emph{Nat.\ Photon.} \textbf{15}, 102--114 (2021).

\bibitem{Silva2014}
A.\,Silva \emph{et al.}, ``Performing mathematical operations with metamaterials,'' \emph{Science} \textbf{343}, 160--163 (2014).

\bibitem{Hughes2019}
T.\,W.\,Hughes \emph{et al.}, ``Wave physics as an analog recurrent neural network,'' \emph{Sci.\ Adv.} \textbf{5}, eaay6946 (2019).

\bibitem{Marcucci2020}
G.\,Marcucci \emph{et al.}, ``The ultranarrowband particle filter: theory, applications, and extensions,'' \emph{Nat.\ Rev.\ Phys.} \textbf{2}, 141 (2020).

\bibitem{Larger2012}
L.\,Larger \emph{et al.}, ``Photonic information processing beyond Turing: an optoelectronic implementation of reservoir computing,'' \emph{Opt.\ Express} \textbf{20}, 3241--3249 (2012).

\bibitem{Tanaka2019}
G.\,Tanaka \emph{et al.}, ``Recent advances in physical reservoir computing,'' \emph{Neural Netw.} \textbf{115}, 100--123 (2019).

\bibitem{Padilla2022}
W.\,J.\,Padilla and R.\,Averitt, ``Metamaterial-enabled wave-based computing,'' \emph{Nat.\ Rev.\ Phys.} \textbf{4}, 507--517 (2022).

\bibitem{Lin2018}
X.\,Lin \emph{et al.}, ``All-optical machine learning using diffractive deep neural networks,'' \emph{Science} \textbf{361}, 1004--1008 (2018).

\bibitem{Davies2018}
M.\,Davies \emph{et al.}, ``Loihi: A neuromorphic manycore processor with on-chip learning,'' \emph{IEEE Micro} \textbf{38}, 82--99 (2018).

\bibitem{Merolla2014}
P.\,A.\,Merolla \emph{et al.}, ``A million spiking-neuron integrated circuit with a scalable communication network and interface,'' \emph{Science} \textbf{345}, 668--673 (2014).

\bibitem{Padamsee2008}
H.\,Padamsee, \emph{RF Superconductivity: Science, Technology, and Applications} (Wiley, 2009).

\bibitem{Romanenko2014}
A.\,Romanenko \emph{et al.}, ``Ultra-high quality factors in superconducting niobium cavities in ambient magnetic fields up to 190\,mG,'' \emph{Appl.\ Phys.\ Lett.} \textbf{105}, 234103 (2014).

\bibitem{Likharev1991}
K.\,K.\,Likharev and V.\,K.\,Semenov, ``RSFQ logic/memory family: A new Josephson-junction technology for sub-terahertz-clock-frequency digital systems,'' \emph{IEEE Trans.\ Appl.\ Supercond.} \textbf{1}, 3--28 (1991).

\bibitem{Keyes1975}
R.\,W.\,Keyes, ``The evolution of digital electronics towards VLSI,'' \emph{IEEE Trans.\ Electron Devices} \textbf{26}, 271--279 (1979).

\bibitem{Lugiato2015}
L.\,A.\,Lugiato, F.\,Prati, and M.\,Brambilla, \emph{Nonlinear Optical Systems} (Cambridge Univ.\ Press, 2015).

\bibitem{Strogatz2018}
S.\,H.\,Strogatz, \emph{Nonlinear Dynamics and Chaos}, 2nd ed.\ (Westview Press, 2015).

\end{thebibliography}

\end{document}
